\documentclass{mcom-l}
% \usepackage[notextcomp,nomath]{stix2}
\usepackage{tcolorbox}
\usepackage{amsfonts,amscd,amssymb,amsmath,amsthm,enumerate}
\numberwithin{equation}{section}
\usepackage{array,mathtools,mathrsfs,bm}
\usepackage{threeparttable}
\usepackage{booktabs,caption,longtable,multicol,multirow}
\usepackage{subcaption,lscape}
\usepackage{makecell}
\usepackage{graphicx}
\usepackage{enumerate}
\usepackage{enumitem}
\setlist[enumerate]{noitemsep, topsep=2pt}
\setlist[itemize]{noitemsep, topsep=2pt}
% \usepackage[linesnumbered,boxruled,lined,commentsnumbered]{algorithm2e}
\usepackage[colorlinks,citecolor=purple,linkcolor=blue]{hyperref}
% \usepackage[nameinlink]{cleveref}
\usepackage{cleveref}
% \usepackage[parfill]{parskip}
\usepackage[margin=1.5in]{geometry}
\usepackage[sort&compress,numbers]{natbib}
% \usepackage{authblk}
\usepackage{algorithm,algorithmic}
\usepackage{pifont}
\usepackage{appendix}
\usepackage{multirow}
\usepackage{tabularx}
\usepackage{makecell}
\usepackage{tikz}
\usepackage{pgfplots}
\usetikzlibrary{arrows}
\usetikzlibrary{calc}
\usetikzlibrary{arrows.meta}
\usetikzlibrary{backgrounds}
\usepgfplotslibrary{patchplots}
\usepgfplotslibrary{fillbetween}
% \usepackage[inline]{showlabels}
% \renewcommand{\showlabelfont}{\tiny\color{blue}}
%%%%%%%%%%%%%%%%%%%%%%%%%%%%%%%%%%%%%%%%%%%%%%%%
\usepackage{eqparbox}
\renewcommand\algorithmiccomment[1]{%
  \hfill\texttt{\#}\ \eqparbox{COMMENT}{\texttt{#1}}%
}
%%%%%%%%%%%%%%%%%%%%%%%%%%%%%%%%%%%%%%%%%%%%%%%%

%% tikz
\definecolor{ao(english)}{rgb}{0.0, 0.5, 0.0}
\definecolor{cadmiumgreen}{rgb}{0.0, 0.42, 0.24}
\definecolor{darkpastelgreen}{rgb}{0.01, 0.75, 0.24}
\let\subfigureautorefname=\figurename
\usepackage{adjustbox}
% Define block styles
\tikzset{%
  materia/.style={draw, fill=blue!20, text width=6.0em, text centered, minimum height=1.5em,drop shadow},
  etape/.style={materia, text width=8em, minimum width=10em, minimum height=3em, rounded corners, drop shadow},
  texto/.style={above, text width=6em, text centered},
  linepart/.style={draw, thick, color=black!50, -LaTeX, dashed},
  line/.style={draw, thick, color=black!50, -LaTeX},
  ur/.style={draw, text centered, minimum height=0.01em},
  back group/.style={fill=white!20,rounded corners, draw=black!50, dashed, inner xsep=15pt, inner ysep=10pt},
}
\pgfplotsset{%
  layers/standard/.define layer set={%
      background,axis background,axis grid,axis ticks,axis lines,axis tick labels,pre main,main,axis descriptions,axis foreground%
    }{
      grid style={/pgfplots/on layer=axis grid},%
      tick style={/pgfplots/on layer=axis ticks},%
      axis line style={/pgfplots/on layer=axis lines},%
      label style={/pgfplots/on layer=axis descriptions},%
      legend style={/pgfplots/on layer=axis descriptions},%
      title style={/pgfplots/on layer=axis descriptions},%
      colorbar style={/pgfplots/on layer=axis descriptions},%
      ticklabel style={/pgfplots/on layer=axis tick labels},%
      axis background@ style={/pgfplots/on layer=axis background},%
      3d box foreground style={/pgfplots/on layer=axis foreground},%
    },
}

\newcommand{\etape}[2]{node (p#1) [etape] {#2}}

\newcommand{\transreceptor}[3]{%
  \path [linepart] (#1.east) -- node [above] {\scriptsize #2} (#3);}
\tikzstyle{matheq} = [node distance=8.75cm, text width=21em, minimum width=1cm,
minimum height=2em, text centered,color=black!50]

\renewcommand{\baselinestretch}{1.2}

%
\newcommand{\gkc}{\|\gk\|^{1/2}}
\newcommand{\xk}{x_k}
\newcommand{\Hk}{H_k}
\newcommand{\gk}{g_k}
\newcommand{\dk}{d_k}
\newcommand{\gkn}{g_{k+1}}
\newcommand{\xkn}{x_{k+1}}
\newcommand{\dkn}{d_{k+1}}
\newcommand{\lk}{\lambda_{k+1}}
\newcommand{\cmark}{\ding{51}}%
\newcommand{\xmark}{\ding{55}}%
\newcommand{\arca}{\texttt{CubicReg-A}}

\newtheorem{theorem}{Theorem}[section]
\newtheorem{lemma}[theorem]{Lemma}

\theoremstyle{definition}
\newtheorem{definition}[theorem]{Definition}
\newtheorem{example}[theorem]{Example}
\newtheorem{xca}[theorem]{Exercise}

\theoremstyle{remark}
\newtheorem{remark}[theorem]{Remark}
\newtheorem{corollary}{Corollary}[section]
% \newtheorem{lemma}{Lemma}[section]
\newtheorem{proposition}{Proposition}[section]
\newtheorem{condition}{Condition}[section]
% \newtheorem{definition}{Definition}[section]
% \newtheorem{example}{Example}[section]
\newtheorem{question}{Question}[section]
% \newtheorem{remark}{Remark}[section]
\newtheorem{assumption}{Assumption}[section]
\newtheorem{strategy}{Strategy}[section]


\def\algorithmautorefname{Algorithm}
\def\subsectionautorefname{Section}
\def\definitionautorefname{Definition}
\def\corollaryautorefname{Corollary}
\def\lemmaautorefname{Lemma}
\def\assumptionautorefname{Assumption}
\def\propositionautorefname{Proposition}
\def\tableautorefname{Table}
\def\strategyautorefname{Strategy}
\def\remarkautorefname{Remark}
\def\conditionautorefname{Condition}


\newcommand{\hsodm}{\texttt{HSODM}}
\newcommand{\hsodmarc}{\texttt{HSODMArC}}
\newcommand{\hsodmhvp}{\texttt{HSODM-HVP}}
\newcommand{\drsom}{\texttt{DRSOM}}
\newcommand{\drsomh}{\texttt{DRSOM-H}}
\newcommand{\lbfgs}{\texttt{LBFGS}}
\newcommand{\newtontr}{\texttt{Newton-TR}}
\newcommand{\cg}{\texttt{CG}}
\newcommand{\arc}{\texttt{ARC}}
\newcommand{\atrms}{\texttt{NO-ATR}}
\newcommand{\itrace}{\texttt{TRACE}}
\newcommand{\gd}{\texttt{GD}}
\newcommand{\utr}{\texttt{UTR}}
\newcommand{\utrhvp}{\texttt{UTR-HVP}}
\newcommand{\newtontrst}{\texttt{Newton-TR-STCG}}
\newcommand{\iutr}{\texttt{iUTR(Hessian)}}
\newcommand{\iutrhvp}{\texttt{iUTR}}

\newcommand{\atr}{\texttt{GL-ATR}}

\makeatletter
\renewcommand\paragraph{\@startsection{paragraph}{4}%
  \z@ \z@ {-.5em}%
  {\normalfont\normalsize\bfseries}}%
\makeatother


\usepackage{float}
\usepackage{etoolbox}
\usepackage{cleveref}

% Define subroutine float
\newfloat{subroutine}{htbp}{los}
\floatname{subroutine}{Subroutine}

% Define cref name for subroutine
\crefname{subroutine}{subroutine}{subroutines}
\Crefname{subroutine}{Subroutine}{Subroutines}

\allowdisplaybreaks
% \usepackage[inline]{showlabels}
%     If you need symbols beyond the basic set, uncomment this command.
%\usepackage{amssymb}

%     If your article includes graphics, uncomment this command.
%\usepackage{graphicx}

%     If the article includes commutative diagrams, ...
%\usepackage[cmtip,all]{xy}


%     Update the information and uncomment if AMS is not the copyright
%     holder.
%\copyrightinfo{2009}{American Mathematical Society}

\numberwithin{equation}{section}

\begin{document}

\title{Our Projects on Second-Order Optimization}

%    Only \author and \address are required; other information is
%    optional.  Remove any unused author tags.

%    author one information
% \author[Jiang et al.]{Yuntian Jiang, Chuwen Zhang, Bo Jiang, Yinyu Ye}
\author{Y. Jiang \ C. He \ C. Zhang \ B. Jiang \ D. Ge \ Y. Ye}
% \address{School of Information Management and Engineering, Shanghai University of Finance and Economics, Shanghai 200433, China}


%    \subjclass is required.
% \subjclass[2020]{Primary 65K05, 90C25, 90C30}

% \date{}

% \dedicatory{}

% \keywords{Trust-Region Method, Acceleration, Global Complexity, Local Analysis}
\maketitle
\section{Universal Trust-Region Method}
We aim to find second-order stationary points~(SOSPs) of the nonconstrained optimization problem:
\begin{equation*}
    \min_{x\in \mathbb{R}^n} f(x).
\end{equation*}
$x$ is a SOSP if
\begin{align*}
     & \|\nabla f(x) \| \leq O(\epsilon)                            \\
     & \lambda_{\min}(\nabla ^2 f(x)) \geq \Omega(-\epsilon^{1/2}).
\end{align*}
\subsection{Motivation}
\begin{enumerate}
    \item The standard trust-region method converges to SOSPs with complexity $O(\epsilon^{-2})$, which is unsatisfactory.
    \item Many variants have been proposed, e.g., \citet{luenberger_linear_2021,curtis2021trust,curtis2017trust,hamad2022consistently}. Their methods are guaranteed to find SOSPs in $O(\epsilon^{-3/2})$ iterations.
    \item It remains unknown if the above variants have improved complexity in the convex setting. Note the cubic Newton method~\citet{nesterov2006cubic} finds $x$ with
          \begin{equation*}
              f(x) - f^* \leq \epsilon
          \end{equation*}
          in $O(\epsilon^{-1/2})$ iterations.
    \item So the question is that can we design a trust-region variant that improves the complexity results in both cases simultaneously?
\end{enumerate}
\subsection{Our Method}
We assume the Hessian of $f$ is $M$-Lipschitz. We propose to use the following trust-region oracle:
\begin{equation}
    \begin{aligned}
        \min_{d \in \mathbb{R}^n} \quad & \frac{1}{2}d^T\left (\nabla^2 f(\xk) + \sigma_k \left \|\nabla f(\xk)\right\|^{1/2} I \right ) d + \nabla f(\xk)^T d \\
        \text{s.t. } \quad              & \|d\| \le \Delta_k := r_k \|\nabla f(\xk)\|^{1/2}.
        \label{eq.newTR}
    \end{aligned}
\end{equation}
Our algorithm is
\begin{algorithm}[!h]
    \caption{A \textbf{U}niversal \textbf{T}rust-\textbf{R}egion Method (UTR)}\label{alg.conceptual utr}
    \begin{algorithmic}[1]
        \STATE \textbf{Input:} Initial point $x_0 \in \mathbb{R}^n$
        \FOR{$k = 0, 1, \ldots, j_f$}
        \STATE Adjust $(\sigma_k, r_k)$ by a proper strategy
        \STATE Solve the subproblem \eqref{eq.newTR} and obtain the direction $d_k$
        \IF{the decrease is sufficient}
        \STATE $x_{k+1} = x_k + d_k$
        \ELSE
        \STATE Go to Line 2
        \ENDIF
        \ENDFOR
    \end{algorithmic}
\end{algorithm}
\subsection{The contribution}
\begin{enumerate}
    \item When $M$ is available, we show that we can always set $(\sigma_k,r_k) = (\frac{\sqrt{M}}{3},\frac{1}{3\sqrt{M}})$. When we apply this strategy, we find first-order stationary points~(FOSPs) in $O(\epsilon^{-3/2})$ iterations in the nonconvex case. Also, when $f$ is convex, we find $\epsilon$-approximate solutions in terms of function value in $O(\epsilon^{-1/2})$ iterations.
    \item When we do not know $M$, we can design an search procedure over $M$. As a result, we devised an adaptive version, which has the same complexity as the fixed version. Differently, the adaptive version guarantees that we find an SOSP.
    \item Compared to other trust-region variants, our assumption is weaker. Except \citet{curtis2021trust}, other methods more or less rely on additional assumptions like $L$-smoothness.
\end{enumerate}
\section{Accelerated Trust-Region Method}
Based on the above method, we take one step further: consider the accelerated variant of the trust-region method in the convex setting.
\subsection{Motivation}
\begin{enumerate}
    \item There are no existing accelerated trust-region-type methods.
    \item Existing accelerated second-order methods have no strong local convergence.
\end{enumerate}
\subsection{Our Method}
We assume $f$ has $M$-Lipschitz Hessian, and $L$-Lipschitz gradient. We use the more general oracle, which no longer tight the parameter to the current gradient norm.
\begin{equation*}
    \begin{aligned}
        \min_{d \in \mathbb{R}^n} \quad m^{\text{TROA}}(d) & := \nabla f(x)^T d + \tfrac{1}{2} d^T (\nabla^2 f(x) + \sigma I) d, \\
        \text{s.t.} \quad \|d\|                            & \le r.
    \end{aligned}
\end{equation*}
We base our algorithm on the outcome of the multiplier $\lambda$:
\begin{enumerate}
    \item When $\lambda>0$, we can directly guarantee the global acceleration.
    \item When $\lambda = 0$, we view it as a hint that the algorithm may enter the local convergence region of Newton's method.
\end{enumerate}
\begin{algorithm}[ht]
    \caption{Global-Local Accelerated Trust-Region Method (\atr{})}\label{alg.accelerated utr}
    \begin{algorithmic}[1]
        \STATE \textbf{input:} initial point $x_0=v_0\in \mathbb{R}^n$, $s_0=0$, tolerance $\epsilon>0$
        \FOR{$k=0, 1, \ldots, K_\epsilon$}
        \STATE $y_k = \frac{k}{k+3}x_k+\frac{3}{k+3}v_k$
        \STATE $(\sigma_k,r_k) = (\tfrac{\sqrt{2M}}{2}\|\nabla f(y_k)\|^{1/2},\tfrac{1}{\sqrt{2M}}\|\nabla f(y_k)\|^{1/2})$
        \STATE $(d_k,\lambda_k) = \operatorname{TROA}(y_k,\sigma_k,r_k)$
        % \noindent\rule{8cm}{0.4pt} 
        \IF{$\lambda_k=0$} \label{line.if}
        % \STATE $x_{k+1}=y_k+d_k$ \label{line.nonzero multiplier}
        \STATE \textbf{\# Local Detection Begins}
        % \STATE $z_k,\text{ET}=\operatorname{Newton}(y_k,\left \lceil\frac{\log \|\nabla f(y_k)\|}{\epsilon} \right\rceil)$ 
        \STATE $(z_k,d_k,\mu_k,\text{ET}) = \operatorname{LD}(y_k,\tfrac{\sqrt{2M}}{2}\|\nabla f(y_k)\|^{1/2},d_k,\epsilon)$
        \label{line.early terminate main}
        \IF{$\text{ET}$}
        \STATE \textbf{terminate and output $z_k$}
        \ENDIF
        % \IF{$\textrm{ET}$} \STATE\textbf{terminate} \ENDIF
        \ENDIF
        \STATE $x_{k+1}=y_k+d_k$
        \STATE $s_{k+1} = s_{k}+\tfrac{(k+1)(k+2)}{2}\nabla f(x_{k+1})$
        \STATE $v_{k+1} = v_0-\sqrt{\tfrac{8}{3M\|s_{k+1}\|}}s_{k+1}$
        \ENDFOR
    \end{algorithmic}
\end{algorithm}
\begin{algorithm}[!ht]
    \caption{Local Detection for Algorithm~\ref{alg.accelerated utr}}\label{alg.local detection}
    \begin{algorithmic}[1]
        \STATE{\textbf{input:} $z_0=y\in \mathbb{R}^n$, $\mu_+$, $d_+$, $\epsilon$;
            \label{line.right bracket point}
        }
        \IF{$\|d_+\| \leq \min\{\frac{\epsilon}{2\kappa_H}, \sqrt{\frac{\epsilon}{M}}\}$}
        \STATE{\textbf{output} $y+d_+$, None, None, ET =True
            \label{line.et 2}
        }
        \ENDIF
        \STATE{\textbf{\# Track 1: Local Region Detection}}
        \FOR{$i=1,\ldots,\left \lceil\log\frac{\|\nabla f(y)\|}{\epsilon} \right\rceil$}
        \STATE{$z_k = z_{k-1}-\nabla^2f(z_{k-1})^{-1}\nabla f(z_{k-1})$}
        \IF{$\|\nabla f(z_k)\|\leq \epsilon$}
        \STATE{\textbf{output} $z_k$, None, None, ET=True
            \label{line.et 1}
        }
        \ENDIF
        \ENDFOR
        \STATE{\textbf{\# Track 2: Ratio Pass Search}}
        \IF{$\frac{\mu_+}{\|d_+\|}\leq 2M$
            \label{line.check right bracket point}}
        \STATE{\textbf{output} None, $d_+$, $\mu_+$, ET =False
            \hfill\texttt{\# Check if the bisection is needed}
        }
        \ELSE
        \STATE{$r_- = \|d_+\|$, $r_+ = \left \|\left(\nabla^2f(y)+M\|d_+\| I\right)^{-1}\nabla f(y)\right\|$
            \label{line.left bracketing point}
        }
        \ENDIF
        % \WHILE{True}
        % \STATE{$r = \tfrac{r_-+r_+}{2}$
        % \hfill\texttt{\# Bisection starts}
        % }
        % \STATE{$(d,\mu) = \operatorname{TROA}(y,0,r)$
        % \label{line.linear system oracle}
        % }
        % \IF{$\frac{\mu}{\|d\|}<M$}
        % \STATE{$r_+=r$}
        % \ELSIF{$\tfrac{\mu}{\|d\|}>2M$}
        % \STATE{$r_- = r$}
        % \ELSE
        % \STATE{\textbf{output} None, $d$, $\mu$, ET= False}
        % \ENDIF
        % \ENDWHILE
        \WHILE{$\frac{\mu}{\|d\|} \not\in [M, 2M]$}
        \STATE{Use bisection method to find $r \in [r_-, r_+]$}
        \hfill\texttt{\# Bisection starts}

        \STATE{$(d,\mu) = \operatorname{TROA}(y,0,r)$
            \label{line.linear system oracle}
        }

        \ENDWHILE
        \STATE{\textbf{output} None, $d$, $\mu$, ET= False}
    \end{algorithmic}
\end{algorithm}
\subsection{Contribution}
\begin{enumerate}
    \item Our first trust-region variant, which combines our new trust-region oracle with Nesterov's acceleration~\citet{nesterov2008accelerating}, shows a global complexity of $\tilde O(\epsilon^{-1/3})$ and a local convergence of $O(\log\log 1/\epsilon)$~(under sufficient second-order condition).
    \item Our second trust-region variant, which combines our oracle with the Monteiro-Svaiter acceleration~\citet{monteiro2013accelerated}, achieves a near-optimal global complexity of $\tilde O(\epsilon^{-2/7})$.
\end{enumerate}
\clearpage
\addcontentsline{toc}{section}{References}

\bibliographystyle{plainnat}
\bibliography{ref}
\end{document}